\begin{frame}{Lectura y representación}
  \begin{columns}[T,onlytextwidth]
    \begin{column}{0.47\textwidth}
      \begin{block}{Instancias usadas}
        \begin{itemize}
          \item BACP8
          \item BACP10
          \item BACP12
        \end{itemize}
      \end{block}
      \begin{exampleblock}{Estado académico}
        \small
        En el semestre actual $t$ se consultan:
        \begin{itemize}
          \item Estado \texttt{passed[i]}: 1 si el curso fue aprobado; 0 si sigue pendiente.
          \item Plan \texttt{assigned\_periods[i]}: semestre del curso $i$.
        \end{itemize}
      \end{exampleblock}
    \end{column}
    \begin{column}{0.49\textwidth}
      \begin{block}{Representación}
        \small
        $A_t$: aprobados; $U_t=C\setminus A_t$: pendientes.\\[1mm]
        \[X_t=(x_i)_{i\in U_t},\qquad x_i\in\{t,\ldots,T\}.\]
        El vector del código tiene $|C|$ posiciones: para $i\in U_t$,
        \texttt{assigned\_periods[i]} guarda $x_i$.\\[2mm]
        Ejemplo: curso $i=2$ y valor $x_i=3$ significa que el curso 2
        quedó planificado para el semestre 3. El 0 indica aprobado o
        pendiente aún sin asignar; \texttt{passed[i]} los distingue.
      \end{block}
      % \begin{alertblock}{Decisión de diseño}
      %   El vector da una sola asignación por curso y permite asociar feromona
      %   a cada par \emph{(curso, semestre)}. Una matriz binaria exigiría
      %   controlar además una asignación única por fila.
      % \end{alertblock}
    \end{column}
  \end{columns}
\end{frame}
